\section{The experiment and its retained information}\label{sec:intro}
An observed history fixes the laws of its executable continuations. These laws identify a predictive quotient, but they do not by themselves identify a finite procedure that can maintain that quotient. The procedure must also select actions, transform reports, retain information between acquisitions and produce a task readout. Its occupation law is generated by those operations, not chosen independently of them.

A particularly consequential distinction occurs at a physical regeneration. A new patient, job, specimen or prepared system may be independent of its predecessor conditional on an unknown fixed model. The observer's finite record need not be independent: it may contain evidence about that model, or a control state that will affect the next experiment. Replacing that record by a prescribed reset label changes the information structure. The correct boundary object is then a finite Markov chain coupled to the durations and rewards of the intervening physical experiments, rather than an iid sequence of complete cycles.

Our main theorem makes this boundary object part of a general experiment-to-risk transfer. It permits arbitrary standard Borel physical state spaces and Borel report rules. It assumes neither a finite-dimensional filter nor a Fisher or covariance geometry. The model parameter is fixed for the entire experiment and is unavailable to the one program being optimized. The raw hypotheses are finite-product response domination and uniform integrability of the physical return time. The additional recurrence quantity is computed from the \emph{actually induced} retained-state transition matrix. It does not count as an extra store of information.

The central mechanism has three parts. First, finite-response continuity is proved for a report row that is reused on distinct acquisition coordinates. Second, a group-inverse bound gives a compact, quantitatively regular class of retained programs even when there are several closed classes or a period. Third, a stopped-cycle argument adds physical overshoot to a finite-state Poisson bias. Thus a return moment and a boundary recurrence scale enter separately:
\begin{equation}\label{eq:intro-rate}
 |J_N-g|\ \le\ {3\bar\mu(2M+1)H\over N}
       +2^{1-p}MK N^{-p},\qquad 0<p\le1.
\end{equation}
The same program has a normalized-discount error bounded by the corresponding expression with $N^{-1}$ and $N^{-p}$ replaced by $1-\beta$ and $(1-\beta)^p$. Neither a controller reset nor a free external phase variable is introduced in these comparisons.

The group inverse and its Poisson interpretation are classical \cite{Meyer}. The finite-model conclusion, once compactness has been established, is a finite-intersection argument rather than a minimax exchange. What is proved here is their joint use for one unknown-model Borel causal program, with retained intercycle state, class-dependent physical time, uniform raw-call comparison and same-memory error certification. The recurrence hypothesis is not hidden: Section~\ref{sec:proof} identifies precisely the compact boundary families on which cost-universal Cesaro convergence is uniform. It also explains how the unrestricted average value of any finite model family is recovered by exhausting these regular classes.

The geometric theorem has a different, explicitly stated role. For a specified exploration, actual class-weighted occupation gives a checkpoint identity and an online lower bound. A stationary, counted controller--predictor lift gives the upper bound. Taking infima yields a joint-control sandwich; additional uniform mass and attainment conditions give a matching optimized memory law. We do not identify an arbitrary fixed-policy quantization rate with an unrestricted control theorem. Likewise, known-law conditional prediction and robust unknown-model control are separate quantifiers unless a common predictive realization has been verified.

Two raw verifications show that the boundary theorem is not a repackaged filtering example. The first is a controlled, partially observed queue with an unbounded workload, where return tails follow from a drift inequality for every legal action. The second uses noncommuting completely positive instruments and heralded termination; the observer retains its classical control state while the physical system is prepared anew. Neither proof is a premise of the general theorem. A worked two-state queue certificate gives the explicit error $17/N$ for an entire feedback class, and a corresponding $17(1-\beta)$ discount bound. The geometric consequences include nonexpansive propagation, actual rank-changing mass and nonuniform length-biased holding laws.

The comparison with long-run POMDP theory is substantive. Finite-memory approximate optimality in \cite{CSZ}, strong uniform values in \cite{VZ}, and average-cost equations in \cite{Borkar} concern broader strategy or model settings but do not give the retained, fixed-cardinality certificate proved here. History-space methods \cite{Yuksel} and finite-state average-cost approximation \cite{YuBertsekas} address complementary existence and approximation questions. Controlled restart \cite{APZ} is not compulsory observer overwrite; intermittent information \cite{ADK} is not free controller memory. Section~\ref{sec:comparison} gives an assumption-by-assumption comparison, including recent complexity results \cite{ACL}. We claim neither a general efficient synthesis algorithm nor a new discovery of finite-memory learning \cite{HellmanCover}.
